\documentclass[10pt]{article}
\usepackage[a4paper,margin=2cm]{geometry}
\usepackage{booktabs,tabularx,array,multirow}
\usepackage[T1]{fontenc}
\usepackage{lmodern}
\newcommand{\TT}{$T\bar T$}
\newcolumntype{Y}{>{\raggedright\arraybackslash}X}

\title{Summary of agent reports on \emph{Matrix theory reloaded} (\texttt{paper.tex})}
\author{}
\date{6 November 2026}

\begin{document}
\maketitle
\vspace{-2em}

\section*{Setup}
Twelve runs: three models (Haiku~5.5, Sonnet~5.5, Opus~5.5) $\times$ two effort levels (low, high) $\times$ two repeats.
Each run wrote a \texttt{report.tex} giving (1)~the principal achievement, (2)~the principal drawback and
(3)~the most worthwhile follow-up. I put each answer into a category based on what the report presents as its
\emph{headline} choice. Points raised only as supporting detail are counted separately in Table~\ref{tab:mentions}.

\section*{1. Principal achievement}
All 12 runs chose the same principal achievement.
\begin{itemize}
\item \textbf{A1 -- Unified BPS-limit framework (12/12).} The BPS decoupling limit that defines matrix theory (M$p$T),
  when applied at the asymptotic infinity of a D$p$-brane geometry, \emph{is} the near-horizon limit. This places matrix theory,
  AdS/CFT and IMSY, non-Lorentzian holography and the DLCQ$^n$/DLCQ$^m$ classification in one scheme.
\end{itemize}
Every run also named the $p$-brane \TT{} flow equations as the main concrete result. Most runs treated them as a
by-product. Two Haiku runs (\texttt{high\_1}, \texttt{low\_2}) gave them equal billing as a second headline
achievement. Two runs checked the flow equations numerically: Haiku \texttt{high\_2} ($p=1,2$) and
Opus \texttt{high\_2} ($p=0,2,3,4,5$, up to 6 scalars). Both found that the equations hold.

\section*{2. Principal drawback}
\begin{itemize}
\item \textbf{D1 -- Kinematic or conjectural only, with no dynamical test (8/12).} The evidence consists of matching limits
  and scalings. No observable (spectrum, correlator, thermodynamics, symmetry algebra) is compared across any new
  duality, so the DLCQ$^n$/DLCQ$^m$ conjecture works as a classification rather than a falsifiable claim.
\item \textbf{D2 -- Unverified non-Lorentzian bulk backgrounds (3/12).} The headline weakness is that the new geometries
  are not shown to solve the M$p$T supergravity equations and torsion constraints, and that they rely on an
  ad hoc rescaling of $\ell$ (that is, of $N$).
\item \textbf{D3 -- The ``intrinsic'' \TT{} generation of the bulk is not derived (1/12).} $\Omega(r)$ is inserted by
  hand from the known solution instead of following from the flow.
\end{itemize}
These categories overlap heavily. Every report mentions the unverified backgrounds somewhere (Table~\ref{tab:mentions}).
The difference lies in which point the report puts first.

\section*{3. Recommended follow-up}
\begin{itemize}
\item \textbf{F1 -- Formulate or use M$p$T supergravity, verify the non-Lorentzian backgrounds and compute their
  isometry and super-algebras (7/12).} Variants differ in the starting case ($p=0$, $p=0,1$, or M2) and in whether
  the supergravity equations must first be derived.
\item \textbf{F2 -- Non-abelian and/or supersymmetric extension of the $p$-brane \TT{} flows (2/12).} Test against the
  Myers action, BFSS/BMN, or BLG/ABJM.
\item \textbf{F3 -- Dynamical holographic test of a specific new duality (2/12).}
  Opus \texttt{high\_1}: DLCQ$^1$/DLCQ$^2$ via the M5 uplift (AdS$_7\times S^4$ with a null circle), matching symmetries, sources and the protected spectrum.
  Opus \texttt{high\_2}: large-$N$ thermodynamics of the $p=2$ (ABJM) flow compared with non-extremal M2-branes in asymptotically flat space.
\item \textbf{F4 -- Derive $\Omega(r)$ from the non-Lorentzian equations of motion (M2 case) (1/12).} This tests the intrinsic-bulk claim.
\end{itemize}

\section*{Breakdown by model and effort}

\begin{table}[h]
\centering
\caption{Headline choice per run. A1 = unified BPS/near-horizon framework. Drawbacks: D1 = no dynamical test,
D2 = unverified backgrounds, D3 = $\Omega(r)$ put in by hand. Follow-ups: F1 = M$p$T supergravity check,
F2 = non-abelian/susy flows, F3 = dynamical holographic test, F4 = derive $\Omega(r)$. The runner-up is the
alternative that the report ranked second, if it named one.}
\label{tab:runs}
\small
\begin{tabularx}{\textwidth}{llccc>{\raggedright\arraybackslash}p{3.4cm}Y}
\toprule
Model & Effort/run & Ach. & Drawback & Follow-up & Follow-up detail & Runner-up \\
\midrule
Haiku  & low\_1  & A1 & D2 & F1 & Start at $p=0$, then $p\le3$ & Susy abelian \TT{} flow \\
Haiku  & low\_2  & A1 & D2 & F2 & Non-abelian flow vs.\ Myers/BFSS/BMN & Supergravity check (F1) \\
Haiku  & high\_1 & A1 & D2 & F1 & CAS check + isometry algebra & Derive $\Omega(r)$ from flow (F4) \\
Haiku  & high\_2 & A1 & D3 & F4 & $p=2$, $\Omega(r)$ unknown in EOM & none ranked \\
\midrule
Sonnet & low\_1  & A1 & D1 & F1 & Formulate EOM, verify, symmetries & Non-abelian/susy flows (F2) \\
Sonnet & low\_2  & A1 & D1 & F2 & M2/D2 flow vs.\ BLG/ABJM, Myers & Supergravity / observables \\
Sonnet & high\_1 & A1 & D1 & F1 & Build sugra ($p=0,1$), then observables & Non-abelian flows (F2) \\
Sonnet & high\_2 & A1 & D1 & F1 & Sugra + isometries + one observable & Non-abelian flows (F2) \\
\midrule
Opus   & low\_1  & A1 & D1 & F1 & Expand sugra ($p=0$ or 2), superalgebra & Non-abelian/susy 3d flow (F2) \\
Opus   & low\_2  & A1 & D1 & F1 & M2T sugra, superalgebra match & Non-abelian/susy flows (F2) \\
Opus   & high\_1 & A1 & D1 & F3 & M0T/D4 via M5 uplift: symmetries, spectrum & Quantise $p=2$ flow \\
Opus   & high\_2 & A1 & D1 & F3 & ABJM $T\bar T$ thermodynamics vs.\ M2 & none ranked \\
\bottomrule
\end{tabularx}
\end{table}

\begin{table}[h]
\centering
\caption{Counts of headline choices by model and effort (two runs per cell).}
\label{tab:counts}
\begin{tabular}{l cc cc cc c}
\toprule
 & \multicolumn{2}{c}{Haiku} & \multicolumn{2}{c}{Sonnet} & \multicolumn{2}{c}{Opus} & \\
\cmidrule(lr){2-3}\cmidrule(lr){4-5}\cmidrule(lr){6-7}
Category & low & high & low & high & low & high & Total \\
\midrule
A1 Unified BPS/near-horizon framework & 2 & 2 & 2 & 2 & 2 & 2 & 12 \\
\midrule
D1 No dynamical test / kinematic only & 0 & 0 & 2 & 2 & 2 & 2 & 8 \\
D2 Unverified non-Lorentzian backgrounds & 2 & 1 & 0 & 0 & 0 & 0 & 3 \\
D3 $\Omega(r)$ inserted by hand & 0 & 1 & 0 & 0 & 0 & 0 & 1 \\
\midrule
F1 M$p$T supergravity check + symmetries & 1 & 1 & 1 & 2 & 2 & 0 & 7 \\
F2 Non-abelian / susy \TT{} flows & 1 & 0 & 1 & 0 & 0 & 0 & 2 \\
F3 Dynamical holographic test & 0 & 0 & 0 & 0 & 0 & 2 & 2 \\
F4 Derive $\Omega(r)$ from EOM & 0 & 1 & 0 & 0 & 0 & 0 & 1 \\
\bottomrule
\end{tabular}
\end{table}

\begin{table}[h]
\centering
\caption{Specific weaknesses mentioned anywhere in the drawback section (not only as the headline).}
\label{tab:mentions}
\begin{tabular}{l cc cc cc c}
\toprule
 & \multicolumn{2}{c}{Haiku} & \multicolumn{2}{c}{Sonnet} & \multicolumn{2}{c}{Opus} & \\
\cmidrule(lr){2-3}\cmidrule(lr){4-5}\cmidrule(lr){6-7}
Weakness & low & high & low & high & low & high & Total \\
\midrule
Backgrounds not shown to solve M$p$T sugra & 2 & 2 & 2 & 2 & 2 & 2 & 12 \\
Flows only classical/bosonic/abelian       & 2 & 2 & 2 & 2 & 2 & 2 & 12 \\
Ad hoc rescaling of $\ell$ (i.e.\ of $N$)  & 2 & 0 & 2 & 2 & 2 & 2 & 10 \\
No observable / quantitative test          & 1 & 1 & 2 & 2 & 2 & 2 & 10 \\
Strong coupling / quantum effects at infinity, $p>3$ & 2 & 1 & 2 & 2 & 2 & 0 & 9 \\
Poincar\'e patch only                      & 1 & 0 & 2 & 1 & 1 & 0 & 5 \\
$\Omega(r)$ put in by hand (circularity)   & 0 & 2 & 0 & 0 & 0 & 2 & 4 \\
Paper excessively long                     & 0 & 0 & 0 & 1 & 0 & 1 & 2 \\
\bottomrule
\end{tabular}
\end{table}

\section*{Observations}
\begin{itemize}
\item \textbf{Achievement:} all runs agree, regardless of model or effort.
\item \textbf{Drawback by model:} every Sonnet and Opus run framed the main weakness as the lack of any \emph{dynamical} test.
  Haiku framed it more narrowly, as specific unverified steps (the backgrounds, or the $\Omega(r)$ input).
\item \textbf{Follow-up:} the modal choice, verifying the non-Lorentzian backgrounds in M$p$T supergravity (F1), is the
  ``safe'' option. It appears in every model. The non-abelian \TT{} extension was the most common runner-up (6 runs) and was the main pick in 2 low-effort runs.
  Only Opus at high effort proposed a dynamical holographic test (F3). Both of those proposals were the most detailed and specific in the set (M5 uplift; ABJM thermodynamics with a $T_{\max}$ comparison).
\item \textbf{Effort:} higher effort produced longer reports (Opus high: 130--166 lines versus 19--35 at low) and more independent checking,
  such as numerical verification of the flow equations. It changed the follow-up choice mainly for Opus.
\end{itemize}

\end{document}
